\section{The Brownian regulator and its transfer operators}
\label{sec:brownian}

The partition comparison used later introduces a model with continuous space
and discrete time.  We first construct its transfer operators and identify the
state on which their norms can be calculated.  The circumference remains
finite throughout.  This point is essential: the comparison with the lattice
model will require a quantitative statement on a specified, growing family
of circles.

\subsection{The row measure and particle expansion}

Let space be the circle $\mathbb R/W\mathbb Z$, where $W>0$, and let one unit
of discrete time be a vertical lattice step.  The spin manifold is
\[
 \mathcal G=(\mathbb R/2\pi\mathbb Z\times SU(2))/
       \bigl((v,Q)\sim(v+\pi,-Q)\bigr),
 \qquad G=e^{iv}Q.
\]
We identify $SU(2)$ with the unit sphere $S^3$.  On a row, use the
unnormalized Brownian loop measure with Hamiltonian
\[
 E=\frac12\bigl(-\partial_v^2-\Delta_{S^3}\bigr),
\]
with Haar probability as reference measure.  Thus the mass of the row
measure is $\operatorname{Tr}e^{-WE}$.  We denote integration with respect to
this measure by $\mathcal E_W$.  For $b>0$, set $\lambda=b^2/2$ and assign
successive rows the oriented interaction
\[
 \exp\left\{\lambda\int_0^W
                \operatorname{Tr}\bigl(G_t(x)G_{t+1}(x)^*\bigr)\,dx\right\}.
\]
Reversing the orientation complex-conjugates the interaction and gives the
same construction after reversing the notation for successive rows.

We use $n_{\mathrm p}$ exclusively for particle number, reserving $N$ for the
renormalization depth later in the paper.  Expand each vertical exponential.
A constant circle rotation on one row multiplies its incoming and outgoing
terms by opposite circle characters.  Integration over that rotation
therefore forces the same particle number on every time-link.  A color
$A_i=(a_i^L,a_i^R)\in\{1,2\}^2$ records the two matrix indices of the $i$th
insertion.  The particle Hilbert space is the bosonic subspace of
\[
 L^2\bigl((\mathbb R/W\mathbb Z)^{n_{\mathrm p}};
                     (\mathbb C^2\otimes\mathbb C^2)^{\otimes n_{\mathrm p}}\bigr).
\]
Equivalently, we integrate positions over the ordered chamber
$0<X_1<\cdots<X_{n_{\mathrm p}}<W$, with the color identifications induced
by simultaneous permutation of particles.  All kernels below act by
integration over this chamber and summation over colors.  The chamber
integral therefore absorbs the factorial in the exponential series.

For position and color lists $(X,A)$ and $(Y,B)$, define
\begin{equation}
 T_{n_{\mathrm p}}(Y,B;X,A)
 =\lambda^{n_{\mathrm p}}\mathcal E_W\left[
       \prod_{i=1}^{n_{\mathrm p}}G_{B_i}(Y_i)
                        \overline{G_{A_i}(X_i)}\right].
 \label{eq:brownian-transfer}
\end{equation}
For time period $n$, the periodic partition function is
\[
 Z^B_+(n,W)=\sum_{n_{\mathrm p}\ge0}
                   \operatorname{Tr}T_{n_{\mathrm p}}^{\,n}.
\]
A minus-center seam in time multiplies each term by $(-1)^{n_{\mathrm p}}$.
The expansion is absolutely convergent at every finite period.  Each
$T_{n_{\mathrm p}}$ is a positive operator: its kernel is the Gram kernel of
the insertion functions in the row loop measure.  Moreover,
\[
 \sum_{n_{\mathrm p}\ge0}\operatorname{Tr}T_{n_{\mathrm p}}
 =e^{2\lambda W}\operatorname{Tr}e^{-WE}<\infty,
\]
because $\sum_{a,b}|G_{ab}|^2=2$.  In particular, for each parity
$\ell\in\{0,1\}$ the maximum
\[
 \tau_\ell=\max_{n_{\mathrm p}\equiv\ell\ ({\rm mod}\ 2)}
                              \|T_{n_{\mathrm p}}\|
\]
is attained.

\begin{proposition}[The parity norm difference]
\label{prop:brownian-parity}
Let $b\to\infty$, and suppose that for a fixed $D>0$
\[
 W\asymp e^{\pi b}b^{-1/2}(\log b)^D.
\]
Suppose also that there is a fixed $C_*>0$ such that every particle number
attaining either parity maximum satisfies
$|n_{\mathrm p}/W-b^2|\le C_*b$.  Then
\begin{equation}
 \log\tau_0-\log\tau_1
       \sim32\sqrt2\,b^{1/2}e^{-\pi b}.
 \label{eq:B1}
\end{equation}
\end{proposition}

The remainder of this section identifies the relevant eigenstate.
Section~\ref{sec:integral} determines the scalar quantity that will give
\eqref{eq:B1}, and Section~\ref{sec:bethe} completes the finite-circle
spectral calculation.  The particle-number hypothesis will be verified
when the Brownian and Wilson partition functions are compared.

\subsection{Positivity and a commuting contact Hamiltonian}

The two copies of $SU(2)$ act on the left and right binary indices.
A fixed-weight sector specifies the number of occurrences of the second
basis vector in each index.  The \emph{balanced sector} has
$\lfloor n_{\mathrm p}/2\rfloor$ such occurrences in each index.  Every
irreducible representation carried by $n_{\mathrm p}$ binary indices
contains this weight, so rotation invariance implies that the norm of
$T_{n_{\mathrm p}}$ occurs in the balanced sector.

For $c>0$, let $P^L_{ij}$ interchange the left indices of particles $i,j$.
On the symmetric Sobolev form domain define
\begin{equation}
 H_c=-\sum_{i=1}^{n_{\mathrm p}}\partial_i^2
       +2c\sum_{i<j}\delta(X_i-X_j)P^L_{ij}.
 \label{eq:B2}
\end{equation}
Here and below the delta interactions are defined by traces on the contact
hyperplanes.  Their infinitesimal form bound with respect to the kinetic
energy gives a closed, bounded-below form.  On a fixed circle it has compact
resolvent.

\begin{lemma}[The state attaining the transfer norm]
\label{lem:brownian-ground}
For fixed $W>0$ and $n_{\mathrm p}\ge0$, the ground state of $H_c$ in every
nonempty fixed-weight sector is simple for each $c>0$.  In a fixed color
sign convention it is strictly positive.  The operator $H_2$ commutes with
$T_{n_{\mathrm p}}$, and the norm of $T_{n_{\mathrm p}}$ is its eigenvalue
on the balanced ground state of $H_2$.
The commutation statement remains valid with corresponding diagonal
boundary twists for the particle sections and the Brownian heat loop.
\end{lemma}

\begin{proof}
For $n_{\mathrm p}=0$ the assertions are immediate.  For positive particle
number, the proof uses a common cone of componentwise nonnegative
wavefunctions.  We first show that the transfer kernel is strictly positive
in this cone and that the contact Hamiltonian has a unique positive ground
state.  The commutation identity will then identify the transfer norm on
that state.

Write the rows of $Q$ as
\[
 Q=\begin{pmatrix}
 c_0e^{i\phi}&s_0e^{i\theta}\\
 -s_0e^{-i\theta}&c_0e^{-i\phi}
 \end{pmatrix},
 \qquad c_0,s_0\ge0,\quad c_0^2+s_0^2=1.
\]
Change the sign of the color $(2,1)$.  Every matrix-entry insertion now
has a nonnegative radial amplitude times a single angular monomial in
$(v,\phi,\theta)$.  In a fixed-weight transfer kernel, the total angular
charge vanishes.  Separation of variables expresses the intervening heat
operators, at each angular frequency, as radial heat operators with
nonnegative killing potential.  Their kernels are strictly positive in
the radial interior.  One may first impose Dirichlet boundaries on compact
subintervals of $0<c_0<1$ and then let these intervals increase.  The two
coordinate degeneracies have codimension two and are polar; the same
separated heat kernel is consequently obtained for every angular
frequency.  This proves almost-everywhere strict positivity of the
fixed-weight kernel of $T_{n_{\mathrm p}}$ in the stated sign convention.

The same convention makes the contact Hamiltonian cooperative.  If either
binary index of a contacting pair agrees, the interaction is diagonal on
the symmetric contact trace.  If both indices differ, the left exchange
changes the number of colors $(2,1)$ by one, and its off-diagonal entry is
negative after the sign change.  To use this observation without imposing
artificial boundaries at coincidences, regard a wavefunction as a list
of species components, each symmetric among particles of the same color.
The diagonal operator in each component is a scalar Laplacian with scalar
contact terms.  Its heat kernel is positivity improving on the full
periodic position torus.

Here is a direct justification that also supplies positivity of contact
traces.  Mollify each scalar delta interaction in its normal difference
coordinate.  For bridge times in $[1,2]$, the expected time-integral of
these mollified interactions is bounded uniformly in the bridge endpoints
and mollifier scale.  Split the time interval in halves and use the free
heat-kernel bounds together with the one-dimensional normal-coordinate
bound $C(1+t^{-1/2})$ near either endpoint.  Feynman--Kac and Jensen's
inequality give a strictly positive lower bound for the mollified scalar
heat kernel.  The hyperplane trace estimate passes this domination to the
contact operator: translation of a trace through normal distance
$\varepsilon$ changes its $L^2$ norm by at most
$C\sqrt\varepsilon\|u\|_{H^1}$, and the associated resolvent forms are
uniformly coercive at a sufficiently negative spectral parameter.
Integration of the heat-kernel bound gives a positive lower bound for the
scalar resolvent applied to any nonzero nonnegative component, including
its contact trace.  A nonnegative source supported only on contact
hyperplanes also has nonnegative resolvent by the form maximum principle.

The absolute-value form inequality permits a ground vector to be chosen
nonnegative.  In its component equations the off-diagonal contacts give
nonnegative sources.  Moving a sufficiently large positive multiple of
the component to the right-hand side produces a nonzero interior source
whenever that component is nonzero; the preceding scalar resolvent
argument then makes it positive, with positive contact trace.  To see that
positivity reaches every species component, write its four
color counts as $n_{11},n_{12},n_{21},n_{22}$.  Their row sums and column
sums are fixed by the two weights.  Thus $n_{11}$ determines all four
counts and ranges over an interval of integers.  The exchange
$(11),(22)\leftrightarrow(12),(21)$ changes $n_{11}$ by one and connects
successive admissible counts.  Positivity therefore propagates throughout
the sector.  Permutations and the positive multiplicities
arising from identical-species symmetry do not affect this argument.
If a real ground vector changed sign, the absolute-value inequality and
the variational principle would make both its positive and negative
parts nonnegative ground vectors with disjoint supports.  This contradicts
the strict positivity just proved.  Ground-state simplicity follows.

It remains to connect this positive state to the Brownian transfer.
The rotor representations are indexed by charge $q\in\mathbb Z$ and spin
$j\in\frac12\mathbb Z_{\ge0}$, subject to $2j\equiv q\pmod2$.
Writing $d=2j+1$, each representation contributes its $d^2$ matrix elements
and has energy
\[
 E_{qj}=\frac{q^2+d^2-1}{2}.
\]
Multiplication by a matrix entry of $G$ raises $q$ by one and changes $d$
by one in either allowed direction.  In the variables
$x=(q+d)/2$, $y=(q-d)/2$, it therefore raises exactly one variable by one.
Let
\[
 P(z)=\frac{4z^3-z}{3},\qquad
 J=P(x)+P(y).
\]
The finite-difference identity $P(z+1)-P(z)=(2z+1)^2$ gives
\[
 [E,[E,G_{ab}]]=[J,G_{ab}],\qquad
 [E,[E,\overline{G_{ab}}]]=-[J,\overline{G_{ab}}].
\]
In any ordering of insertion positions, derivatives of the rotor trace
insert these commutators.  Cyclicity makes the sum of the $J$ commutators
zero.  Thus the sums of second derivatives in output and input positions
agree away from coincidences.

At a coincidence of two output insertions, the derivative jump in either
coordinate is
\[
 [[E,G_{ab}],G_{cd}]=2G_{cb}G_{ad}.
\]
This follows from the Pauli identity, including the scalar generator;
equivalently, for $E=-\Delta/2$ the double commutator is the negative
scalar product of the two gradients.  The two kinetic derivatives hence
give the contact term canceled by $2cP^L\delta$ at $c=2$.
The conjugate computation handles two inputs, and at a mixed input-output
coincidence the two delta terms cancel.  We have proved the distributional
intertwining identity for $H_2$ and the kernel
\eqref{eq:brownian-transfer}.

All differentiations can be justified before taking distributions.  In
each ordering, some heat interval has length at least $W/(2n_{\mathrm p})$;
multiplication by an entry shifts each angular label by only a bounded
amount.  The rotor series therefore converges with its one-sided
derivatives; in particular, first derivatives are uniformly bounded on
each closed ordering chamber.  Continuity
across its faces makes the kernel globally Sobolev $H^1$ in either set
of positions; intersections of several faces create no further boundary
terms.  The derivative bounds and the finite volume imply that
$T_{n_{\mathrm p}}:L^2\to H^1$ is bounded.  If $\mathfrak h_2$ is the form
of $H_2$, integration by parts in the distributional identity gives
\[
 \mathfrak h_2(T_{n_{\mathrm p}}u,v)
       =\mathfrak h_2(u,T_{n_{\mathrm p}}v)
\]
first for smooth symmetric functions and then, by the trace bound and
form-domain density, for all $u,v\in H^1$.  For $u\in D(H_2)$ the right
side equals $\langle T_{n_{\mathrm p}}H_2u,v\rangle$, by self-adjointness
of the transfer.  Thus $T_{n_{\mathrm p}}u\in D(H_2)$ and
$H_2T_{n_{\mathrm p}}u=T_{n_{\mathrm p}}H_2u$, which also proves
resolvent commutation.

At zero twist, commutation maps the simple balanced ground space of $H_2$
to itself.  Its positive ground vector is therefore an eigenvector of
$T_{n_{\mathrm p}}$.  Strict positivity of the latter kernel identifies
that eigenvalue with its spectral radius in the sector, hence with its
norm.  The weight argument preceding the lemma gives the full norm.

It remains to extend commutation to the matching boundary twists used
later.  For real $\zeta,\chi$, put
$g=\operatorname{diag}(e^{i\chi},e^{-i\chi})$ and
$\varepsilon_2=\left(\begin{smallmatrix}0&1\\-1&0\end{smallmatrix}\right)$.
After the fixed change of exterior color basis
$\widetilde G=\varepsilon_2G\varepsilon_2^T$, use the rotor pullback
\[
 (R_{\zeta,\chi}F)(v,Q)=F(v+\zeta,g^{-1}Qg^{-1}).
\]
Since $\varepsilon_2g^{-1}=g\varepsilon_2$, this pullback transforms
$\widetilde G$ into $e^{i\zeta}g\widetilde Gg$.
Moving an output insertion through the seam therefore gives the particle
color holonomy $A=e^{i\zeta}g\otimes g$; an input insertion has the
conjugate covariance.  The transfer consequently maps sections with
$\psi(\ldots,X_i+W,\ldots)=A_i\psi(\ldots,X_i,\ldots)$ to sections with
the same boundary condition.  On rotor harmonics of charge $q$ and spin
$j$, $R_{\zeta,\chi}$ acts by $e^{i\zeta q}$ and $g_j^{-1}$ on both
matrix indices.  It commutes with $E$ and $J$, so cyclicity and the contact
calculation remain valid.  For these real twists the loop measure is
positive and the boundary holonomy is unitary; the same Gram and weak-form
arguments apply on the twisted Sobolev domain.  Write
$\sigma_3=\operatorname{diag}(1,-1)$, with superscripts denoting its action
on the indicated binary index.  Multiplication by
\[
 \prod_i\exp\!\left\{\frac{iX_i}{W}
       \bigl(\zeta I+\chi(\sigma_3^L+\sigma_3^R)_i\bigr)\right\}
\]
identifies this domain with the periodic one.  In that fixed gauge the
operator identities extend holomorphically to small complex twists;
self-adjointness and positivity are used only for real twists.
\end{proof}

\section{The integral equations and the exponentially small mass}
\label{sec:integral}

We next determine the scalar quantities needed to construct and evaluate
the ground state of $H_2$.  The finite-circle calculation in
Section~\ref{sec:bethe} will organize $n_{\mathrm p}=2M+\ell$ particles
into $M$ pairs with momenta close to $u_j-i,u_j+i$, where $u_j$ is real,
and, when $\ell=1$, one unpaired particle with momentum close to zero.
The density of pair centers per unit circumference will be described by
$\rho_B$ on an interval $(-B,B)$.  Thus $M/W$ is approximated by
$\int_{-B}^B\rho_B$.

The same calculation will express the logarithm of the transfer norm as
a sum of one contribution for each pair and one for the unpaired particle,
with controlled errors.  The equations below reorganize this sum into a common bulk term and a
positive cost $m_*$ for odd particle number.  Their endpoint condition
selects the interval needed to maximize over particle numbers; the precise
finite-circle argument is given in \eqref{eq:B22} and
Lemma~\ref{lem:bethe-parity-error}.  We solve these equations first because
their density and endpoint estimates are also needed in the construction
of those finite-circle states.

The leading exponential of $m_*$ is insufficient for the exact mass:
a correction of order one in $B$ changes its multiplicative constant.
The main analytic task is therefore to match the interior density to its
endpoint profile with an error that retains this correction.

\subsection{Scalar equations and the selected endpoint}

All momenta in this section use the normalization $c=2$ in
\eqref{eq:B2}.  Define
\[
 K_a(x)=\frac{a}{\pi(a^2+x^2)},\qquad K=K_4,
 \qquad S(x)=\frac14\operatorname{sech}(\pi x/2).
\]
Convolution is on $\mathbb R$.  For an interval $(-B,B)$, a density defined
only on that interval is extended by zero unless stated otherwise.
Consider
\begin{equation}
 \begin{aligned}
 \rho_B&=\frac1\pi+K*\rho_B&&(|x|<B),\\
 \epsilon&=p+K*\epsilon_{\rm in},&
 p(x)&=4\log b-\log\bigl((x^2+1)(x^2+9)\bigr),\qquad \epsilon(B)=0.
 \end{aligned}
 \label{eq:B3}
\end{equation}
Here the second equation defines $\epsilon$ on the whole line,
$\epsilon_{\rm in}=\epsilon\mathbf1_{(-B,B)}$, and the boundary condition
selects $B=B_*(b)$.  The function $p$ is the leading logarithmic transfer
contribution associated with the momenta $x-i,x+i$:
\[
 \frac{b^4}{((x-i)^2+4)((x+i)^2+4)}
 =\frac{b^4}{(x^2+1)(x^2+9)}=e^{p(x)}.
\]
We call $\epsilon$ the dressed energy.  Its defining equation is chosen
to cancel the interaction kernel when a sum of pair weights is replaced
by a density integral.  Indeed, symmetry of $K$ and the two equations
in \eqref{eq:B3} give
\[
 \int_{-B}^B p(x)\rho_B(x)\,dx
 =\int_{-B}^B\epsilon(x)\bigl(\rho_B-K*\rho_B\bigr)(x)\,dx
 =\frac1\pi\int_{-B}^B\epsilon(x)\,dx.
\]
The next lemma proves that $\epsilon$ is positive precisely on the selected
interval.  This sign will allow the finite-circle maximization to fill
that interval and to bound all other choices by the same bulk term.

\begin{lemma}[Parametrization by the endpoint]
\label{lem:brownian-equations}
For every $B>0$, the first equation in \eqref{eq:B3} has a unique positive
solution.  There is a unique value $b=b(B)>\sqrt3$ for which the second
has an even solution satisfying $\epsilon(B)=0$; this solution is positive
inside $(-B,B)$ and negative outside.  The function $b(B)$ is continuous
and strictly increasing.  If $\eta_B=-\epsilon'$ on $(-B,B)$, then
\[
 (1-K*)\eta_B=
      \frac{2x}{x^2+1}+\frac{2x}{x^2+9},\qquad |x|<B,
\]
and
\begin{equation}
 4\log b=\log9+\int_0^B a(x)\eta_B(x)\,dx,
 \qquad a(x)=\frac2\pi\arctan\frac4x.
 \label{eq:B4}
\end{equation}
The inverse of $1-K*$ on $(-B,B)$ has supremum norm $O(1+B)$.
\end{lemma}

\begin{proof}
The kernel is positive, and its mass on $(-B,B)$ is bounded above by
$1-c/(1+B)$ uniformly in its center in the interval.  A Neumann series
therefore gives the inverse and the norm bound, as well as positivity of
$\rho_B$.  On odd functions, restriction to $(0,B)$ replaces the kernel
by $K(x-y)-K(x+y)$, which is strictly positive for $x,y>0$.
Iteration gives a unique odd solution $\eta_B$, positive for $x>0$.
Increasing $B$ increases this solution on the old positive half-interval.

Define the even function $\epsilon$ inside the interval by
$\epsilon(x)=\int_{|x|}^B\eta_B(y)\,dy$.
Its boundary values vanish.  Differentiating $K*\epsilon_{\rm in}$
introduces no endpoint terms, so the equation for $\eta_B$ shows that
$(1-K*)\epsilon$ has the required logarithmic form up to an additive
constant.  At zero, changing the order of integration gives
\[
 \epsilon(0)-(K*\epsilon_{\rm in})(0)
 =\int_0^B\left(1-2\int_0^xK(y)\,dy\right)\eta_B(x)\,dx
 =\int_0^B a(x)\eta_B(x)\,dx.
\]
This is \eqref{eq:B4}.  It proves strict monotonicity and continuity of
$b(B)$, and $b(0)=\sqrt3$ by continuation.  Extend $\epsilon$ by the
second equation in \eqref{eq:B3}.  To the right of $B$,
\[
 -\epsilon'(x)=\frac{2x}{x^2+1}+\frac{2x}{x^2+9}
                      +(K*\eta_B)(x)>0,
\]
where positivity follows by pairing the odd integrand at $y$ and $-y$.
Thus the extension is negative outside, as claimed.
\end{proof}

With Fourier transform $\widehat u(k)=\int e^{-ikx}u(x)\,dx$, one has
$\widehat K_a=e^{-a|k|}$ and $\widehat S=(2\cosh k)^{-1}$.  Consequently
\[
 S*(\delta_0-K)=K_1-K_3,
 \qquad
 f(x):=\log\frac{b^2}{x^2+4}=(S*p)(x).
\]
For the logarithmic identity, differentiate first and then fix the
constant using growth at infinity.  Define $\epsilon_{\rm out}
=\epsilon-\epsilon_{\rm in}$ and
\begin{equation}
 m_*=-\left\{f(0)-((K_1-K_3)*\epsilon_{\rm in})(0)\right\}
       =-(S*\epsilon_{\rm out})(0)>0.
 \label{eq:brownian-mass}
\end{equation}
Here $f(0)$ is the leading logarithmic transfer contribution of an
unpaired particle at zero momentum.  As shown in Lemma~\ref{lem:bethe-centers}, that particle
also adds the source $-W^{-1}(K_1-K_3)$ to the pair counting equation.
On a fixed interval, the density approximation gives a logarithmic
correction equal to $W$ times the pairing of this density change with $p$.
Transferring the symmetric resolvent
to $p$ therefore gives the correction
$-((K_1-K_3)*\epsilon_{\rm in})(0)$.  Thus the expression in braces
is the net odd contribution.  The second expression in \eqref{eq:brownian-mass}
follows by convolving \eqref{eq:B3} with $S$ and using the preceding
kernel identity.  It both proves positivity, by the exterior sign of
$\epsilon$, and isolates the part to be evaluated: the exponentially
weighted tail beyond the selected endpoint.

\begin{proposition}[Endpoint, density, and mass asymptotics]
\label{prop:brownian-asymptotics}
As $B\to\infty$, with $b=b(B)$ and $m_*$ defined above,
\begin{equation}
 \begin{aligned}
 b(B)&=\frac B2+\frac{\log(4\pi B)+1}{\pi}+o(1),\\
 m_*&\sim\frac{8}{\pi e\sqrt B}\,e^{-\pi B/2},\\
 2\int_{-B}^B\rho_B(x)\,dx
       &=\frac{B^2}{4}+\frac B\pi\log B+O(B).
 \end{aligned}
 \label{eq:B5}
\end{equation}
There are positive constants $c_1,C_1$ such that, for all sufficiently
large $B$, $\rho_B\ge c_1\sqrt B$ within distance one of either endpoint,
and $\rho_B\le C_1B$ throughout the interval.  For every $B>0$, the
solution also satisfies $\sup_{y\in\mathbb R}|\epsilon'(y)|\le\pi/2$.
\end{proposition}

We prove the proposition in three stages: an outer solution on the scaled
interval, a half-line solution resolving its endpoints, and the final
integrations.  The uniform matching estimate between the first two stages
is what preserves the multiplicative constant in $m_*$.

\subsection{The scaled equations and their outer solutions}

The whole-line odd solution of the equation for $\eta_B$ is
\[
 \eta_\infty(x)=\frac\pi2\tanh(\pi x/2),
\]
as follows by Fourier transformation.  Comparison with this solution in the
positive odd-kernel equation gives $0\le\eta_B(x)\le\pi/2$ for $0<x<B$.
In the next calculation $x$
denotes scaled position.  Set $\delta=4/B$ and define on $\mathbb R$
\[
 h(x)=\eta_\infty(Bx)-\eta_B(Bx),
 \qquad w(x)=\pi\delta\rho_B(Bx),
\]
where $\eta_B$ and $\rho_B$ have been extended by zero.  Their equations are
\begin{equation}
 (1-K_\delta*)h=0,\qquad
 (1-K_\delta*)w=\delta\qquad (|x|<1).
 \label{eq:brownian-scaled}
\end{equation}
The exterior data are $\eta_\infty(Bx)$ and zero respectively; the former
differ uniformly exponentially little from $(\pi/2)\operatorname{sgn}x$.
The function $h$ is odd and nonnegative on the positive half-line, by the
comparison just noted.  Hence its Cauchy convolution is nonnegative there.
Using the whole-line equation gives
$-\epsilon'(y)=\eta_\infty(y)-(K_\delta*h)(y/B)\le\pi/2$ for $y>0$;
the positivity already proved gives the other inequality.  This establishes
the derivative bound in Proposition~\ref{prop:brownian-asymptotics}.

The Cauchy kernel is the Poisson kernel for the upper half-plane.  This
lets us construct an interior approximation by analytic functions.  For
the functions below, sublinear growth at infinity and
integrable boundary singularities give the Poisson representation: if
$u=\operatorname{Re}(F_0+\delta F_1)$ on the real line, then
\[
 (K_\delta*u)(x)=\operatorname{Re}(F_0+\delta F_1)(x+i\delta).
\]
Away from the endpoints, Taylor expansion therefore gives
\[
 (1-K_\delta*)u
 =\delta\operatorname{Im}F_0'
  +\delta^2\left(\tfrac12\operatorname{Re}F_0''
                         +\operatorname{Im}F_1'\right)
  +\text{higher-order terms}.
\]
To solve \eqref{eq:brownian-scaled} to this order, the first coefficient
must be zero for $h$ and one for $w$, and the second must vanish.
The real boundary values must also have the prescribed step and zero
exterior data.

Continue
\[
 r(z)=\sqrt{1-z^2},\qquad
 L(z)=\log\frac{1+z}{1-z}
\]
from $(-1,1)$ through the upper half-plane, with $r$ positive on the
interval.  The functions $\arcsin z$ and $r(z)+iz$ satisfy the leading
conditions for $h$ and $w$, respectively.  The term $iz$ cancels the
linear growth of $r$ at infinity, as required for its Poisson
representation, and supplies $\operatorname{Im}F_{0,w}'=1$.
The following first corrections cancel the second coefficient:
\begin{equation}
 \begin{array}{c|cc}
   &F_0(z)&F_1(z)\\ \hline
 h&\arcsin z&\displaystyle
                \frac{(L(z)-i\pi)/(2\pi)+C_0z}{r(z)}\\[6pt]
 w&r(z)+iz&\displaystyle
                \frac{-z(L(z)-i\pi)/(2\pi)-C_0}{r(z)}.
 \end{array}
 \label{eq:B6}
\end{equation}
Here $C_0$ is a real constant to be determined by endpoint matching.
Its terms have real boundary values on the interval and purely imaginary
boundary values outside, so the interior derivative conditions and the
exterior data do not determine it.  We write $F_{j,h}$ and $F_{j,w}$ for
the two rows.  Direct differentiation gives
\[
 \operatorname{Im}F_{0,h}'=0,\qquad
 \operatorname{Im}F_{0,w}'=1,\qquad
 \operatorname{Im}F_1'=-\tfrac12\operatorname{Re}F_0''.
\]
Thus the table supplies the required interior expansion for every $C_0$.
Its inverse-square-root singularities prevent uniform control at the
endpoints.  An exact half-line solution will both remove those
singularities and fix $C_0$.

\subsection{The half-line solution and uniform endpoint matching}

At the right endpoint the leading interior terms satisfy
\[
 \frac\pi2-\operatorname{Re}\arcsin x\sim\sqrt{2(1-x)},
 \qquad r(x)\sim\sqrt{2(1-x)}.
\]
Put $X=(1-x)/\delta$.  On this scale the kernel becomes $K_1$, and
both $\pi/2-h$ and $w$, divided by $\sqrt{2\delta}$, must match a
profile growing as $\sqrt X$.  To leading order their exterior values
at this edge are zero, and the forcing $\delta$ in the equation for $w$
vanishes after division by $\sqrt{2\delta}$.  We therefore solve the
homogeneous convolution equation on $X>0$ with zero exterior value and
this normalization at infinity.
The profile on the opposite side of the edge will determine the tail
in \eqref{eq:brownian-mass}.

For a related finite-interval Wiener--Hopf analysis of the delta-function
Fermi gas, see Tracy and Widom~\cite{TracyWidom2016}.  We derive the
factorization and the endpoint error estimates needed for the present
kernel.

\begin{lemma}[Half-line solution]\label{brownian:halfline}
There are functions $U,V:[0,\infty)\to\mathbb R$, continuous from the
right at zero and continuous on $(0,\infty)$, with the following
properties.  Extend $U$ by zero on $(-\infty,0)$.  Then
\[
 U(X)=(K_1*U)(X)\quad(X>0),
 \qquad V(Y)=(K_1*U)(-Y)\quad(Y\ge0),
\]
and $U(X)>0$, $V(Y)>0$ for $X,Y\ge0$.  For $\operatorname{Re}s>0$,
\begin{equation}\label{eq:B7}
\begin{split}
 C(s)&=\frac{\sqrt{s}\,
  \exp\!\left\{\frac{s}{2\pi}
             \left(\log\frac{s}{2\pi}-1\right)\right\}}
 {\Gamma(1+s/(2\pi))},\\
 \int_0^\infty e^{-sX}U(X)\,dX
   &=\frac{\sqrt\pi}{2sC(s)},\qquad
 \int_0^\infty e^{-sY}V(Y)\,dY
   =\frac{\sqrt\pi}{2s}C(s).
\end{split}
\end{equation}
The logarithm and square root in this formula are the principal branches.
Moreover,
\begin{equation}\label{brownian:U-asymptotic}
\begin{aligned}
 U(X)&=\sqrt X+
       \frac{\log X+C'}{4\pi\sqrt X}
       +O\!\left(X^{-3/2}(1+\log^2X)\right),\\
 C'&=\log(8\pi)+1,\qquad X\longrightarrow\infty,
\end{aligned}
\end{equation}
while $V(Y)=O((1+Y)^{-1/2})$ and
$U(0)=V(0)=\sqrt\pi/2$.
\end{lemma}

\begin{proof}
We construct $U$ and $V$ from the two transforms in \eqref{eq:B7}, and
then verify the convolution equation.  Both transforms are analytic in
$\operatorname{Re}s>0$.  Stirling's formula gives, uniformly at infinity
in the closed right half-plane,
\[
 C(s)=1+O(s^{-1}),\qquad
 \frac{\sqrt\pi}{2sC(s)}=
 \frac{\sqrt\pi}{2s}+O(s^{-2}),\qquad
 \frac{\sqrt\pi}{2s}C(s)=
 \frac{\sqrt\pi}{2s}+O(s^{-2}).
\]
The Bromwich inverses therefore define locally bounded functions
supported on $[0,\infty)$.  More explicitly, subtract
$\sqrt\pi/(2(s+1))$ from either transform.  The remainder is absolutely
integrable on each vertical line in the right half-plane, so its inverse
is continuous; closing the contour to the right shows that it vanishes
for negative arguments.  The subtracted transform has inverse
$(\sqrt\pi/2)e^{-X}$ on $X>0$.  The $O(s^{-2})$ remainder also shows
that its inverse tends to zero as $X\downarrow0$.  This proves local
boundedness, continuity for positive arguments, and the asserted
right limits at zero.
Both transforms commute with complex conjugation, so their inverses
are real-valued.

To find the large-$X$ behavior, expand at the origin.  With
$\gamma_E$ denoting Euler's constant,
\[
 C(s)=\sqrt{s}\left[
 1+\frac{s}{2\pi}
       \bigl(\log s-\log(2\pi)-1+\gamma_E\bigr)
 +O\!\left(s^2(1+|\log s|^2)\right)\right].
\]
Consequently the transform of $U$ is
\[
 \frac{\sqrt\pi}{2}s^{-3/2}
 -\frac{\sqrt\pi}{4\pi}s^{-1/2}
       \bigl(\log s-\log(2\pi)-1+\gamma_E\bigr)
 +O\!\left(|s|^{1/2}(1+|\log s|^2)\right).
\]
The first term inverts to $\sqrt X$.  Differentiating
$\mathcal L(X^{\alpha-1})(s)=\Gamma(\alpha)s^{-\alpha}$ at
$\alpha=1/2$, and using
$\Gamma'(1/2)/\Gamma(1/2)=-\gamma_E-2\log2$, gives the second term
of \eqref{brownian:U-asymptotic}, including $C'=\log(8\pi)+1$.
For the remainder, deform the inversion contour to the rays
$\arg s=\pm3\pi/4$, joined by a small arc of radius $1/X$.
The principal branches are analytic in the intervening region, and
Stirling's formula is uniform in this sector.  On the small part of
the contour the displayed remainder contributes
$O(X^{-3/2}(1+\log^2X))$ after the substitution $s=t/X$;
on the remaining rays the exponential $e^{sX}$ supplies the required
decay.  This justifies the termwise inversion and its error estimate.
The same contour argument applied to the transform of $V$, whose
leading term at zero is $(\sqrt\pi/2)s^{-1/2}$, gives
$V(Y)=O(Y^{-1/2})$ for large $Y$.  Together with local boundedness,
this is the stated bound on $V$.

It remains to prove that the inverses satisfy the half-line equation.
For real $k\ne0$, the gamma-function identity
\[
 \left|\Gamma\left(1+\frac{ik}{2\pi}\right)\right|^2
  =\frac{|k|/2}{\sinh(|k|/2)}
\]
gives
\begin{equation}\label{brownian:factorization}
 C(ik)C(-ik)=1-e^{-|k|}.
\end{equation}
Use the Fourier convention $\widehat f(k)=\int e^{-ikX}f(X)\,dX$.
The polynomial growth just proved allows distributional Fourier
transforms; the Laplace transforms give their boundary values away
from $k=0$.  Since $\widehat K_1(k)=e^{-|k|}$,
\eqref{brownian:factorization} implies
\[
 (1-e^{-|k|})\frac{\sqrt\pi}{2ikC(ik)}
   =-\frac{\sqrt\pi\,C(-ik)}{2(-ik)}
 \qquad(k\ne0).
\]
Thus the Fourier transform of
$U-K_1*U+\widetilde V$ is supported at zero, where
$\widetilde V(X)=V(-X)$ for $X<0$ and $\widetilde V(X)=0$ for $X>0$.
This assertion can equivalently be checked after inserting a positive
Laplace regulator and then testing against functions supported away
from zero.  A distribution whose Fourier transform is supported at
one point is a polynomial, so $U-K_1*U+\widetilde V$ is a polynomial.

That polynomial vanishes.  Indeed, write
$U(X)=\sqrt{X_+}+R(X)$, where $X_+=\max(X,0)$.
By \eqref{brownian:U-asymptotic}, $R$ is bounded and tends to zero at
both infinities, and hence $K_1*R$ also tends to zero there.
The Poisson extension identity
\[
 (K_1*\sqrt{\,\cdot\,_+})(X)=\operatorname{Re}\sqrt{X+i}
\]
shows that $\sqrt{X_+}-K_1*\sqrt{\,\cdot\,_+}$ tends to zero at
both infinities as well.  Finally $\widetilde V$ tends to zero.
Therefore the polynomial tends to zero and is identically zero.
Restricting this identity to $X>0$ and $X<0$ proves the two
convolution identities.  Their values at zero follow by continuity
of $K_1*U$ and the already established one-sided limits.

For positivity, $U$ is positive for all sufficiently large $X$ by
\eqref{brownian:U-asymptotic}, and its right limit at zero is positive.
If its minimum on the positive half-line were negative, it would be
attained at some $X>0$.  Averaging against $K_1$, with the zero
exterior value, would give a value strictly larger than that minimum,
contradicting $U=K_1*U$.  Thus $U\ge0$; the same averaging equation
and positivity of the kernel make the inequality strict.  The formula
for $V$ then gives $V>0$.
\end{proof}

The normalization $U(X)\sim\sqrt X$ matches the leading outer terms.
Matching the next term of \eqref{brownian:U-asymptotic} to
\eqref{eq:B6} at the right endpoint requires
\[
 \log2+2\pi C_0=\log\delta-C'.
\]
We therefore fix
\[
 C_0=\frac{\log\delta-\log(16\pi)-1}{2\pi}.
\]
The left endpoint gives the same value by symmetry.  We now replace the
two singular outer expansions by this exact profile, subtracting their
common terms so they are counted only once.  Define
\[
 X_+(x)=\frac{1-x}{\delta},\qquad
 X_-(x)=\frac{1+x}{\delta},
\]
and
\[
 D_0(X)=
 \begin{cases}
 U(X)-\displaystyle\left[\sqrt X+
              \frac{\log X+C'}{4\pi\sqrt X}\right],&X>0,\\[6pt]
 0,&X<0.
 \end{cases}
\]
The singularity of $D_0$ at zero is integrable.  Its decay at infinity
follows from \eqref{brownian:U-asymptotic}; in particular,
$D_0\in L^1(\mathbb R)$.  In formulas containing an endpoint,
we use the limit from the side under consideration.

\begin{lemma}[Uniform endpoint matching]\label{brownian:uniform}
Let $h,w$ solve the scaled interval equations stated above, with
$B=4/\delta$, and let $F_{j,h},F_{j,w}$ denote the two rows of
\eqref{eq:B6}.  Define, on the real line away from the endpoints,
\begin{align*}
 h_{\mathrm{app}}(x)
  &=\operatorname{Re}\bigl(F_{0,h}(x)+\delta F_{1,h}(x)\bigr)
    +\sqrt{2\delta}\bigl[-D_0(X_+(x))+D_0(X_-(x))\bigr],\\
 w_{\mathrm{app}}(x)
  &=\operatorname{Re}\bigl(F_{0,w}(x)+\delta F_{1,w}(x)\bigr)
    +\sqrt{2\delta}\bigl[D_0(X_+(x))+D_0(X_-(x))\bigr].
\end{align*}
There exist constants $C<\infty$ and $\delta_0>0$ such that, for
$0<\delta<\delta_0$,
\[
 \|h-h_{\mathrm{app}}\|_{L^\infty(-1,1)}
 +\|w-w_{\mathrm{app}}\|_{L^\infty(-1,1)}
 \le C\delta^{3/2}|\log\delta|.
\]
Both approximants have finite one-sided limits at the endpoints.
\end{lemma}

\begin{proof}
We first verify the cancellation that makes the approximants bounded.
At the right endpoint set $d=1-z$.  The singular terms of the $h$
row of \eqref{eq:B6} are
\[
\begin{aligned}
 F_{0,h}(z)&=\frac\pi2-\sqrt{2d}+O(d^{3/2}),\\
 \delta F_{1,h}(z)&=
 -\frac{\delta}{2\pi\sqrt{2d}}
       \bigl(\log(d/\delta)+C'+i\pi\bigr)\\
 &\quad+O\!\left(\delta\sqrt{|d|}
       (1+|\log\delta|+|\log|d||)\right).
\end{aligned}
\]
Here the constant in the leading correction is exactly $C'$ because
$2\pi C_0=\log\delta-\log(16\pi)-1$.
The $w$ row has the opposite square-root and inverse-square-root
terms, in addition to its analytic terms.  At the left endpoint the
corresponding formulas are obtained by reflection and complex
conjugation.  They give precisely the signs used in the definitions
of $h_{\mathrm{app}}$ and $w_{\mathrm{app}}$.

For clarity, the analytic expression for the two terms subtracted
from $U$ is
\begin{equation}\label{brownian:edge-analytic}
 A(X)=\sqrt X+
       \frac{\log X+C'+i\pi}{4\pi\sqrt X},
 \qquad \operatorname{Im}X<0.
\end{equation}
Its real boundary value is the bracket defining $D_0$ for $X>0$,
and is zero for $X<0$.  Since $X=(1-z)/\delta$ lies in the lower
half-plane when $z$ lies in the upper half-plane, this is the
appropriate continuation at the right edge; at the left edge we use
its conjugate.  Thus adding the $D_0$ terms replaces the displayed
singular terms by multiples of the locally bounded function $U$.
This proves the assertion about endpoint limits.

We estimate the residual before applying the interval comparison
principle.  Each function $U(X_\pm(x))$ satisfies
$(1-K_\delta*)U(X_\pm)=0$ for $-1<x<1$, by the scaling of the
half-line equation.  After separating these exact terms from either
approximant, write its remaining part as
$\operatorname{Re}(G_0+\delta G_1)$.  Here $G_0,G_1$ are obtained
from the corresponding row of \eqref{eq:B6} by subtracting the
analytic edge expressions in \eqref{brownian:edge-analytic}, with
the signs just verified.  The subtraction preserves the identities
\[
 \operatorname{Im}G_0'=q,\qquad
 \operatorname{Im}G_1'=-\tfrac12\operatorname{Re}G_0'',
 \qquad
 q=\begin{cases}0&\text{for }h,\\1&\text{for }w.\end{cases}
\]
In particular $\operatorname{Im}G_0'''=0$ on the interval.
The removed terms leave endpoint remainders of order $d^{3/2}$
in $G_0$ and
$\sqrt d(1+|\log\delta|+|\log d|)$ in $G_1$.

The remaining analytic functions are bounded at the endpoints and
have sublinear growth at infinity, so convolution with $K_\delta$
is their Poisson extension:
$(K_\delta*\operatorname{Re}G)(x)=\operatorname{Re}G(x+i\delta)$.
Put $d_0=1-|x|$.  If $d_0\ge2\delta$, expand $G_0$ to third order
and $G_1$ to first order on the vertical segment from $x$ to
$x+i\delta$.  The derivative identities cancel the term of order
$\delta^2$, and the cubic term from $G_0$ has zero real part.
The derivative bounds supplied by the endpoint remainders therefore
give, for either approximant $f_{\mathrm{app}}$,
\begin{equation}\label{brownian:matching-residual}
 \left|(1-K_\delta*)f_{\mathrm{app}}(x)-q\delta\right|
 \le C\left[
 \delta^4d_0^{-5/2}
 +\delta^3(1+|\log\delta|+|\log d_0|)d_0^{-3/2}
 \right].
\end{equation}
If $d_0\le2\delta$, use the endpoint remainders directly at
$x$ and $x+i\delta$.  Their contributions are bounded by
$C\delta^{3/2}|\log\delta|$.  The derivative of the nonanalytic
$O(d^{3/2})$ part of $G_0$ tends to zero at the endpoint.  Hence
$\operatorname{Im}G_0'=q$ identifies the imaginary part of the analytic
linear coefficient there as $q$, also at the left endpoint by reflection.
The analytic terms therefore contribute
$q\delta+O(\delta^2(1+|\log\delta|))$.
Hence the residual in this region is also at most
$C\delta^{3/2}|\log\delta|$.

On the exterior, the correction from the nearest endpoint vanishes.
The argument of $D_0$ at the opposite endpoint is at least
$2/\delta$, so \eqref{brownian:U-asymptotic} bounds that correction
by $C\delta^2(1+\log^2\delta)$.  The uncorrected outer functions
have exactly the step and zero exterior data, respectively.
The step differs from the exterior data of $h$ by an exponentially
small quantity.  Thus the exterior discrepancy of either approximant
is bounded by
$e_\delta=C\delta^2(1+\log^2\delta)$.

Finally, the mass of the kernel that leaves the interval is
\[
 \lambda_\delta(x)
  =\int_{\mathbb R\setminus(-1,1)}K_\delta(x-y)\,dy
  \ge c\frac{\delta}{d_0+\delta}.
\]
Dividing \eqref{brownian:matching-residual} by this lower bound,
and using the bound in the region $d_0\le2\delta$, gives
\[
 \left|(1-K_\delta*)f_{\mathrm{app}}(x)-q\delta\right|
 \le C\delta^{3/2}|\log\delta|\,\lambda_\delta(x)
 \qquad(-1<x<1).
\]
For example, when $d_0\ge2\delta$, the two ratios are bounded by
$C\delta^3d_0^{-3/2}$ and
$C\delta^2(1+|\log\delta|+|\log d_0|)d_0^{-1/2}$,
respectively, and both have the asserted size.
The nonnegative function
\[
 e_\delta+M\delta^{3/2}|\log\delta|\,
                \mathbf 1_{(-1,1)}
\]
is consequently a comparison bound for both signs of the difference
between the exact solution and its approximant, once $M$ is large
enough.  Indeed the constant has zero Poisson difference, while the
interval indicator has Poisson difference $\lambda_\delta$ inside.
The positivity of $K_\delta$ and the strict loss of mass on the
interval give the comparison principle.  Since
$e_\delta=o(\delta^{3/2}|\log\delta|)$, this proves the claimed
uniform estimate.
\end{proof}

\subsection{Recovering the endpoint, density, and mass}

We now finish the proof of Proposition~\ref{prop:brownian-asymptotics}.
First integrate the approximation for $w$.  Since
$\int_{-1}^1r(x)\,dx=\pi/2$, $\int_{-1}^1r(x)^{-1}\,dx=\pi$, and
$\int_{-1}^1xL(x)/r(x)\,dx$ is finite, its outer part gives
\[
 \int_{-1}^1w(x)\,dx
       =\frac\pi2+\delta\bigl(-\pi C_0+O(1)\bigr).
\]
The matching terms are $O(\delta^{3/2})$, because $D_0$ is integrable, and
the uniform error contributes $O(\delta^{3/2}|\log\delta|)=o(\delta)$.
Returning to the unscaled density gives
\[
 2\int_{-B}^B\rho_B
 =\frac{2B}{\pi\delta}\int_{-1}^1w
 =\frac{B^2}{4}-2BC_0+O(B)
 =\frac{B^2}{4}+\frac B\pi\log B+O(B).
\]
In a physical distance one from either edge, the argument of $U$ ranges
over $[0,1/4]$.  Its positive minimum there, together with the uniform
matching error, gives $w\ge c\sqrt\delta$ and hence
$\rho_B\ge c_1\sqrt B$.  The same composite approximation is bounded
throughout the interval, which gives $\rho_B\le C_1B$.

To recover $b(B)$, use \eqref{eq:B4} and write
$\eta_B(y)=\eta_\infty(y)-h(y/B)$.  The whole-line contribution is
\begin{equation}
 \int_0^B a(y)\eta_\infty(y)\,dy
             =4\log B-\log9+o(B^{-1}).
 \label{eq:brownian-wholeline-integral}
\end{equation}
Indeed, $a=(1-K*)\operatorname{sgn}$ on the positive half-line.
Transfer $1-K*$ from the sign function to $\eta_\infty$ in the
symmetrized integral over $(-B,B)$.  Write
$\eta_\infty=(\pi/2)\operatorname{sgn}+r_\infty$, where $r_\infty$ is
odd and exponentially decreasing.  The sign part cancels in the
crossing-end error.  Oddness of $r_\infty$ and the difference of the two
Cauchy kernels make the remaining error $O(B^{-2})$.
The remaining integral is
\[
 \int_0^B\left(\frac{2y}{y^2+1}+\frac{2y}{y^2+9}\right)\,dy
 =\log\frac{(B^2+1)(B^2+9)}9,
\]
which proves \eqref{eq:brownian-wholeline-integral}.

For the leading outer term $H_0(x)=\operatorname{Re}\arcsin x$,
\begin{equation}
 \int_0^B a(y)H_0(y/B)\,dy
                   =4\log2-\frac8B+o(B^{-1}).
 \label{eq:brownian-leading-integral}
\end{equation}
The constant and first correction follow respectively from
\[
 \frac8\pi\int_0^1\frac{H_0(x)}x\,dx=4\log2,
 \qquad
 \int_0^\infty y\left(a(y)-\frac8{\pi y}\right)\,dy=-8.
\]
For the second identity, integrate with respect to the scale in the
Cauchy kernel.  To justify its use in
\eqref{eq:brownian-leading-integral}, note that
$B H_0(y/B)\to y$ and $B H_0(y/B)\le C y$ for $0<y<B$;
the displayed correction is absolutely integrable.

For the real part of $\delta F_{1,h}$ one similarly obtains
\begin{equation}
 \int_0^B a(y)\delta\operatorname{Re}F_{1,h}(y/B)\,dy
                 =\frac{8+16C_0}{B}+o(B^{-1}).
 \label{eq:brownian-correction-integral}
\end{equation}
Here the required elementary identity is
\[
 \int_0^1\frac{L(x)}{x r(x)}\,dx=\frac{\pi^2}{2}.
\]
For example, replace $L(x)$ by
$\log((1+\ell x)/(1-\ell x))$, differentiate in $\ell$, and integrate
$0\le\ell\le1$.  After replacing $a(Bx)$ by $8/(\pi Bx)$, the two
terms in $F_{1,h}$ contribute $8/B$ and $16C_0/B$.
The replacement error is $o(B^{-1})$: at zero the numerator is linear
in $x$, and its coefficient grows only as $1+|C_0|=O(\log B)$.

The edge compensation in the $h$ approximation contributes $o(B^{-1})$
to this integral.  Near the right edge this is bounded by
$CB^{-3/2}\int_0^\infty|D_0(X)|\,dX$; the term coming from the opposite
edge is smaller by the decay of $D_0$.  The uniform approximation error
contributes at most
\[
 C B^{-3/2}\log B\int_0^B a(y)\,dy
           =O(B^{-3/2}\log^2B)=o(B^{-1}).
\]
Combining \eqref{eq:brownian-wholeline-integral}--
\eqref{eq:brownian-correction-integral} in \eqref{eq:B4} therefore gives
\[
 4\log b=4\log(B/2)-\frac{16C_0}{B}+o(B^{-1}).
\]
Since $2\pi C_0=-\log(4\pi B)-1$, exponentiation proves the first
line of \eqref{eq:B5}.  In particular, $b(B)\to\infty$, so the endpoint
$B_*(b)$ is defined for every sufficiently large $b$.

It remains to evaluate the exponentially small quantity
\eqref{eq:brownian-mass}.  This uses the endpoint profile on the exterior,
rather than an integral of the interior approximation.  The whole-line
equation for $\eta_\infty$ gives, also for $y>B$,
\[
 -\epsilon'(y)=\eta_\infty(y)-(K_\delta*h)(y/B).
\]
At $y=B+4Y$, the uniform endpoint approximation and Poisson convolution
imply
\begin{equation}
 \frac{-\epsilon'(B+4Y)}{\sqrt{2\delta}}\longrightarrow V(Y),
 \qquad Y\ge0,
 \label{eq:brownian-exterior-profile}
\end{equation}
with a bound independent of $Y\ge0$ and sufficiently large $B$.
To see both assertions, write
\[
 \frac{\pi/2-h(1-\delta X)}{\sqrt{2\delta}}.
\]
On $X\ge0$ this converges to $U(X)$ and is bounded by
$C\sqrt{1+X}$.  For $0\le X\le(2\delta)^{-1}$ the bound follows from
the local composite approximation.  For $X\ge(2\delta)^{-1}$, including
the opposite endpoint and its exterior, use $|h|\le\pi/2$ to bound the
quotient by $C\delta^{-1/2}\le C'\sqrt X$.  On $X<0$ its value is
exponentially small.
Convolution at the exterior point is with $K_1(X+Y)$.  The bound is
integrable against that kernel, uniformly in $Y\ge0$, and hence gives
\eqref{eq:brownian-exterior-profile} and the stated domination.

Since $\epsilon$ is even and vanishes at the endpoints,
\[
 m_*=2\int_B^\infty S(y)
                      \int_B^y[-\epsilon'(z)]\,dz\,dy.
\]
Use $S(y)\sim\frac12e^{-\pi y/2}$, interchange the integrals, and put
$z=B+4Y$.  The domination just proved permits passage to the limit and
gives
\begin{equation}
 m_*\sim16e^{-\pi B/2}\sqrt{2\delta}\,
                \frac1{2\pi}\int_0^\infty e^{-2\pi Y}V(Y)\,dY.
 \label{eq:brownian-mass-evaluation}
\end{equation}
Formula~\eqref{eq:B7} yields
\[
 C(2\pi)=\frac{\sqrt{2\pi}}e,
 \qquad
 \int_0^\infty e^{-2\pi Y}V(Y)\,dY=\frac1{2\sqrt2\,e}.
\]
Substitution of $\delta=4/B$ into
\eqref{eq:brownian-mass-evaluation} proves the second line of
\eqref{eq:B5}, and completes Proposition~\ref{prop:brownian-asymptotics}.

For later use, the conversion between the two large parameters is
\[
 e^{-\pi B/2}=4\pi B e\,e^{-\pi b}(1+o(1)),
 \qquad
 m_*\sim32\sqrt2\,b^{1/2}e^{-\pi b}.
\]
Thus it will suffice to identify the parity difference with $m_*$ to
error $O(W^{-1})$: under the width assumption of
Proposition~\ref{prop:brownian-parity}, $W^{-1}=o(m_*)$.
The density estimate also gives $2\int\rho_{B_*}=b^2+O(b)$, matching the
particle-number scale in that proposition.
